
本章考察紧致\BookRecordIndex{T0358}辛流形 $(M,\omega)$ 中 \BookRecordIndex{T0003}$J$-全纯曲线的\BookRecordIndex{T0246}模空间，并证明：在适当假设下，
对通有的\BookRecordIndex{T0088}殆复结构 $J$，这个空间是光滑有限维流形。我们将考虑定义在复结构
$j_\Sigma$ 固定的紧黎曼曲面 $\Sigma$ 上的 \BookRecordIndex{T0003}$J$-全纯曲线。

区分两种情形十分重要。若要对所有（非常值）\BookRecordIndex{T0003}$J$-全纯曲线实现\BookRecordIndex{T0166}横截性，就必须或者允许
\BookRecordIndex{T0088}殆复结构族由 $\Sigma$ 参数化，或者允许适当的\BookRecordIndex{T0157}哈密顿扰动。若\BookRecordIndex{T0088}殆复结构与
$\Sigma$ 中的基点无关，那么只能对单纯 \BookRecordIndex{T0003}$J$-全纯曲线实现\BookRecordIndex{T0166}横截性。本章限于\BookRecordIndex{T0091}单纯曲线的
情形；更一般的情形将在第 6.7 节讨论。

第~\ref{sec:3-1}节引入\BookRecordIndex{T0091}单纯曲线的\BookRecordIndex{T0246}模空间，并讨论这种情形下的主要结果。
第~\ref{sec:3-2}节处理
Thom--Smale \BookRecordIndex{T0166}横截性，并为\BookRecordIndex{T0091}单纯曲线和通有\BookRecordIndex{T0088}殆复结构建立相应的\BookRecordIndex{T0166}横截性定理。由此可知，
只要每条 \BookRecordIndex{T0003}$J$-全纯曲线的\BookRecordIndex{T0346}线性化算子都是满射，\BookRecordIndex{T0246}模空间就是光滑流形。我们还将证明，
\BookRecordIndex{T0246}模空间带有自然定向。第~\ref{sec:3-3}节描述四维中的一些简单例子；在那里，
对所有\BookRecordIndex{T0088}殆复结构都可
直接证明\BookRecordIndex{T0346}线性化算子为满射。在这一特殊情形，论证依赖 Riemann--Roch 定理
（见附录 C），而无需 萨德--斯梅尔定理。

初次阅读时可以略去最后两节。第~\ref{sec:3-4}节讨论满足\BookRecordIndex{T0415}逐点约束的曲线的\BookRecordIndex{T0246}模空间；
这与第 6 章中\BookRecordIndex{T0340}稳定映射\BookRecordIndex{T0246}模空间的\BookRecordIndex{T0166}横截性讨论有关。第~\ref{sec:3-5}节证明适用于当前背景的
\BookRecordIndex{T0380}隐函数定理的一个定量
版本：在适当假设下，它保证\BookRecordIndex{T0193}近似 $J$-全纯曲线附近存在 \BookRecordIndex{T0003}$J$-全纯曲线。该结果被表述成
\BookRecordIndex{T0254}粘合论证所需的形式，并且对 $\Sigma$ 上一类适当的体积形式给出一致估计。

虽然我们没有显式展开满足拉格朗日边界条件的曲线的性质，但确实为理解这种情形奠定了
许多基础。相关讨论和应用见第~\ref{sec:9-2} 节；尤其参见注~\ref{rem:9-2-3}。

\section{\BookRecordIndex{T0091}单纯曲线的\BookRecordIndex{T0246}模空间}\label{sec:3-1}
\zhindex{danchunquxianmofangjian}{单纯曲线的模空间}
\zhindex{hengjiexing}{横截性}

本节陈述单纯 \BookRecordIndex{T0003}$J$-全纯曲线情形下的主要结果\sourcefootnote{原书此处写作
“\textit{simple $J$-holomorphic spheres}”；但本节随即固定任意紧黎曼曲面
$\Sigma$，且定理~\ref{thm:3-1-6} 与定理~\ref{thm:3-1-8} 也在这一范围内陈述，
故改作“曲线”。}，证明则在第~\ref{sec:3-2}节给出。这两节展示了主要\BookRecordIndex{T0408}正则性论证的结构。
\BookRecordIndex{T0091}单纯曲线的性质已在第~\ref{sec:2-5}节讨论。

设 $(M,\omega)$ 是紧致辛 $2n$-流形。固定一个紧黎曼曲面
$(\Sigma,j_\Sigma,d\operatorname{vol}_\Sigma)$ 和一个
$\omega$-\BookRecordIndex{T0372}驯服\BookRecordIndex{T0088}殆复结构 $J\in\mathcal J_\tau(M,\omega)$。我们考虑方程
\begin{equation}\label{eq:3-1-1}
  \bar\partial_J(u)=0,
\end{equation}
其中
\[
  \bar\partial_J(u):=\frac12\bigl(du+J\circ du\circ j_\Sigma\bigr).
\]
给定同调类 $A\in H_2(M;\mathbb Z)$，用
\[
  \mathcal M(A,\Sigma;J)
  :=\bigl\{u\in C^\infty(\Sigma,M)\bigm|
  J\circ du=du\circ j_\Sigma,\ [u]=A\bigr\}
\]
表示 \eqref{eq:3-1-1} 中代表类 $A$ 的解的\BookRecordIndex{T0246}模空间。其单纯解子空间记作
\[
  \mathcal M^*(A,\Sigma;J)
  :=\bigl\{u\in\mathcal M(A,\Sigma;J)\bigm|u\text{ 是单纯的}\bigr\}.
\]
在 $\Sigma=\mathbb{CP}^1$ 的情形，简记
\[
  \mathcal M(A;J):=\mathcal M(A,\mathbb{CP}^1;J),
  \qquad
  \mathcal M^*(A;J):=\mathcal M^*(A,\mathbb{CP}^1;J).
\]

\BookRecordIndex{T0246}模空间 $\mathcal M(A,\Sigma;J)$ 可以解释为一个无穷维向量丛的截面的零点集，具体如下。
令 $\mathcal B\subset C^\infty(\Sigma,M)$ 表示所有代表同调类
$A\in H_2(M;\mathbb Z)$ 的光滑映射 $u:\Sigma\to M$ 所成的空间。可把这个空间看作
一个无穷维（Fréchet）流形；它在 $u\in\mathcal B$ 处的切空间为
\[
  T_u\mathcal B=\Omega^0(\Sigma,u^*TM),
\]
即沿 $u$ 的所有光滑向量场 $\xi(z)\in T_{u(z)}M$ 所成的空间。考虑无穷维向量丛
$\mathcal E\to\mathcal B$，其在 $u$ 处的纤维为
\[
  \mathcal E_u=\Omega^{0,1}(\Sigma,u^*TM),
\]
即 $\Sigma$ 上取值于 $u^*TM$ 的光滑 $J$-反线性 $1$-形式所成的空间。于是 $du$ 的
复反线性部分定义了这个向量丛的一个截面
$\mathcal S:\mathcal B\to\mathcal E$：
\[
  \mathcal S(u):=(u,\bar\partial_J(u)).
\]
\BookRecordIndex{T0246}模空间 $\mathcal M(A,\Sigma;J)$ 是这个截面的零点集。子集
$\mathcal M^*(A,\Sigma;J)$ 则是 $\mathcal M(A,\Sigma;J)$ 与开子集
$\mathcal B^*\subset\mathcal B$ 的交，其中 $\mathcal B^*$ 由所有代表类 $A$ 且\BookRecordIndex{T0397}在某处
单射的光滑映射 $u:\Sigma\to M$ 组成。

\paragraph{算子 \texorpdfstring{$D_u$}{Du}。}
本书的目标是说明如何为一类适当的\BookRecordIndex{T0358}辛流形构造\BookRecordIndex{T0143}格罗莫夫--威滕不变量。这个计划的第一步
是证明：对通有的\BookRecordIndex{T0088}殆复结构 $J$，\BookRecordIndex{T0246}模空间 $\mathcal M^*(A,\Sigma;J)$ 是具有适当维数的
光滑有限维流形。为此必须证明 $\mathcal S$ 与零截面横截。这意味着：对每个
$u\in\mathcal M^*(A,\Sigma;J)$，微分
$d\mathcal S(u):T_u\mathcal B\to T_{(u,0)}\mathcal E$ 的像都与零截面的切空间
$T_u\mathcal B$ 互补。给定 $u\in\mathcal M^*(A,\Sigma;J)$，记
\[
  D_u:=D_{J,u}:=D\mathcal S(u):
  \Omega^0(\Sigma,u^*TM)\longrightarrow\Omega^{0,1}(\Sigma,u^*TM)
\]
为微分 $d\mathcal S(u):T_u\mathcal B\to T_{(u,0)}\mathcal E$ 与投影
\[
  \pi_u:T_{(u,0)}\mathcal E=T_u\mathcal B\oplus\mathcal E_u
  \longrightarrow\mathcal E_u
\]
的复合。算子 $D_u$ 称为截面 $\mathcal S$ 在 $u$ 处的\emph{\BookRecordIndex{T0310}竖直微分}。于是\BookRecordIndex{T0166}横截性可以
表述为：对每个 $u\in\mathcal M^*(A,\Sigma;J)$，\BookRecordIndex{T0346}线性化算子 $D_u$ 都是满射。

在 $\Sigma$ 与 $M$ 的局部坐标中，一条 \BookRecordIndex{T0003}$J$-全纯曲线是方程
\[
  \partial_su+J(u)\partial_tu=0
\]
的解 $u:\mathbb C\to\mathbb R^{2n}$，而沿 $u$ 的向量场是映射
$\xi:\mathbb C\to\mathbb R^{2n}$。沿 $\xi$ 方向对该方程左端求导，便得到
$D_u\xi$ 的局部公式：
\[
  D_u\xi=\eta\,ds-J(u)\eta\,dt,
  \qquad
  \eta:=\frac12\bigl(\partial_s\xi+J(u)\partial_t\xi
  +\partial_\xi J(u)\partial_tu\bigr).
\]
利用 $u$ 是 $J$-全纯的这一事实，可得
\begin{equation}\label{eq:3-1-2}
  D_u\xi=\bar\partial_J\xi
  -\frac12\bigl(J\partial_\xi J\bigr)(u)\partial_J(u).
\end{equation}
注意，右端第二项是反全纯 $1$-形式，因为
\[
  \partial_J(u):=\frac12\bigl(du-J\circ du\circ j_\Sigma\bigr)
\]
是全纯 $1$-形式，而 $J$ 与 $J(\partial_\xi J)$ 反对易。以上讨论表明，由公式
\eqref{eq:3-1-2} 定义的算子与局部坐标的选取无关。直接验证这一点留给读者。

由公式 \eqref{eq:3-1-2}，$D_u$ 是附录 C 意义下的实线性\BookRecordIndex{T0210}柯西--黎曼算子，因而是
\BookRecordIndex{T0128}弗雷德霍姆算子。有关 $D_u$ 的这些信息足以建立本章大部分结果。然而，第 3.5 节中的
\BookRecordIndex{T0380}隐函数定理与第 10 章中的\BookRecordIndex{T0256}粘合定理依赖精细估计，而这些估计建立在对算子 $D_u$ 的透彻
理解之上。因此，值得花些功夫来发展 $D_u$ 的简洁公式。

$D_u$ 的定义可以推广到一般的光滑映射 $u:\Sigma\to M$，无论它是否 $J$-全纯；
但此时定义依赖于把切空间 $T_{(u,\bar\partial_J(u))}\mathcal E$ 分裂为水平子空间和
竖直子空间的一种选取。这种分裂依赖于 $TM$ 上的联络。我们必须假设该联络保持\BookRecordIndex{T0088}殆复结构
$J$，从而 $\mathcal E$ 的纤维在逐点\BookRecordIndex{T0263}平行移动下不变。\footnote{设
$\mathbb R\to\mathcal B:\lambda\mapsto u_\lambda$ 是一条光滑曲线。若对每个
$\zeta\in T_z\Sigma$，沿曲线 $\lambda\mapsto u_\lambda(z)$ 的向量场
$\lambda\mapsto\alpha_\lambda(z;\zeta)\in T_{u_\lambda(z)}M$ 都是平行的，则称
$\mathcal E$ 沿此曲线的截面
$\alpha_\lambda\in\Omega^{0,1}(\Sigma,u_\lambda^*TM)$ 关于 $TM$ 上的联络
$\widetilde\nabla$ 是平行的。若 $\xi(\lambda)\in T_{x(\lambda)}M$ 是沿 $M$ 中一条曲线
的任意向量场，则相应的协变导数记作 $\widetilde\nabla_\lambda\xi$。}
我们采用由度量 \eqref{eq:2-1-1} 的 \BookRecordIndex{T0034}Levi--Civita 联络 $\nabla$ 诱导的复线性联络
\[
  \widetilde\nabla_vX
  :=\nabla_vX-\frac12J(\nabla_vJ)X.
\]

给定 $\xi\in\Omega^0(\Sigma,u^*TM)$，令
\[
  \Phi_u(\xi):u^*TM\longrightarrow\exp_u(\xi)^*TM
\]
表示如下复向量丛同构：沿测地线
$s\mapsto\exp_{u(z)}(s\xi(z))$ 关于 $\widetilde\nabla$ 作\BookRecordIndex{T0263}平行移动。
（这里和往常一样，测地线总是关于 $\nabla$ 取的。）然后定义映射
\[
  \mathcal F_u:\Omega^0(\Sigma,u^*TM)
  \longrightarrow\Omega^{0,1}(\Sigma,u^*TM)
\]
为
\begin{equation}\label{eq:3-1-3}
  \mathcal F_u(\xi)
  :=\Phi_u(\xi)^{-1}\bar\partial_J\bigl(\exp_u(\xi)\bigr).
\end{equation}
这个映射正是截面 $\mathcal S:\mathcal B\to\mathcal E$ 在由
$\widetilde\nabla$ 确定的平凡化下的竖直部分。因此，对每条 \BookRecordIndex{T0003}$J$-全纯曲线
$u:\Sigma\to M$，$\mathcal F_u$ 在零点的微分就是上面定义的算子 $D_u$。对一般映射
$u$，我们将以这个导数作为 $D_u$ 的定义。下面的命题给出一个显式公式。

\begin{proposition}\label{prop:3-1-1}
对任意光滑映射 $u:\Sigma\to M$，定义算子
\[
  D_u:\Omega^0(\Sigma,u^*TM)\longrightarrow\Omega^{0,1}(\Sigma,u^*TM)
\]
为
\[
  D_u\xi:=d\mathcal F_u(0)\xi.
\]
则对每个 $\xi\in\Omega^0(\Sigma,u^*TM)$，都有
\begin{equation}\label{eq:3-1-4}
  D_u\xi
  =\frac12\bigl(\nabla\xi+J(u)\nabla\xi\circ j_\Sigma\bigr)
  -\frac12J(u)(\nabla_\xi J)(u)\partial_J(u).
\end{equation}
\end{proposition}

\begin{proof}
考虑道路
\[
  \mathbb R\longrightarrow C^\infty(\Sigma,M):
  \lambda\longmapsto u_\lambda:=\exp_u(\lambda\xi).
\]
依定义，$\Phi_u(\lambda\xi)\mathcal F_u(\lambda\xi)=\bar\partial_J(u_\lambda)$。
由于 $\Phi_u(\lambda\xi)$ 由沿测地线 $\lambda\mapsto u_\lambda(z)$ 的\BookRecordIndex{T0263}平行移动给出，
所以
\[
\begin{aligned}
  D_u\xi
  &=
  \left.\frac{d}{d\lambda}\mathcal F_u(\lambda\xi)\right|_{\lambda=0}\\
  &=
  \left.\widetilde\nabla_\lambda\bar\partial_J(u_\lambda)
  \right|_{\lambda=0}\\
  &=
  \left.\frac12\bigl(
    \widetilde\nabla_\lambda du_\lambda
    +J(u_\lambda)\widetilde\nabla_\lambda du_\lambda\circ j_\Sigma
  \bigr)\right|_{\lambda=0}\\
  &=
  \left.\frac12\bigl(
    \nabla_\lambda du_\lambda
    +J(u_\lambda)\nabla_\lambda du_\lambda\circ j_\Sigma
  \bigr)\right|_{\lambda=0}\\
  &\quad
  -\left.\frac14\bigl(
    J(u_\lambda)(\nabla_{\partial_\lambda u_\lambda}J)(u_\lambda)du_\lambda
    -(\nabla_{\partial_\lambda u_\lambda}J)(u_\lambda)du_\lambda\circ j_\Sigma
  \bigr)\right|_{\lambda=0}\\
  &=
  \frac12\bigl(\nabla\xi+J(u)\nabla\xi\circ j_\Sigma\bigr)
  -\frac12J(u)(\nabla_\xi J)(u)\partial_J(u).
\end{aligned}
\]
最后一个等式用到了 $\nabla$ 无挠这一事实。它蕴含：对每个光滑映射
$\mathbb R^2\to M:(\lambda,t)\mapsto\gamma(\lambda,t)$，都有
$\nabla_\lambda\partial_t\gamma=\nabla_t\partial_\lambda\gamma$；
特别地，这适用于映射 $\gamma(\lambda,t)=u_\lambda(z(t))$。
\end{proof}

\begin{remark}[自然性]\label{rem:3-1-2}
有时在记号中指出算子 $D_u$ 对复结构 $j=j_\Sigma$ 的依赖较为方便，此时写作
$D_{j,u}:=D_u$。于是映射
\[
  (j,u,\xi)\longmapsto D_{j,u}\xi
\]
在 $\Sigma$ 的微分同胚群作用下等变，即
\begin{equation}\label{eq:3-1-5}
  D_{\phi^*j,\phi^*u}\phi^*\xi=\phi^*(D_{j,u}\xi)
\end{equation}
对每个光滑映射 $u:\Sigma\to M$、每个沿 $u$ 的向量场
$\xi\in\Omega^0(\Sigma,u^*TM)$、$\Sigma$ 上每个（殆）复结构 $j$，以及每个微分同胚
$\phi:\Sigma\to\Sigma$ 都成立。这里记
\[
  \phi^*u:=u\circ\phi,\qquad
  \phi^*\xi:=\xi\circ\phi,\qquad
  \phi^*\alpha:=\alpha\circ d\phi
\]
（其中 $\alpha\in\Omega^{0,1}(\Sigma,u^*TM)$）。特别地，如果复结构 $j$ 固定而
$\phi$ 是 $\Sigma$ 的全纯微分同胚，那么
\[
  D_{\phi^*u}\phi^*\xi=\phi^*(D_u\xi).
\]
\end{remark}

\medskip\noindent\textbf{习题。}
证明
$\mathcal F_{\phi^*j,\phi^*u}(\phi^*\xi)=\phi^*\mathcal F_{j,u}(\xi)$，
并由此\BookRecordIndex{T0326}推出公式 \eqref{eq:3-1-5}。证明在 $M$ 的微分同胚群作用下的自然性。

\begin{remark}\label{rem:3-1-3}
若 $J$ 与 $\omega$ \BookRecordIndex{T0348}相容，则算子 $D_u$ 还可写成
\begin{equation}\label{eq:3-1-6}
  D_u\xi=(\widetilde\nabla\xi)^{0,1}
  +\frac14N_J\bigl(\xi,\partial_J(u)\bigr),
\end{equation}
其中 $N_J$ 表示 $J$ 的 \BookRecordIndex{T0042}Nijenhuis 张量，而
$(\widetilde\nabla\xi)^{0,1}$ 表示 $1$-形式
$\widetilde\nabla\xi\in\Omega^1(\Sigma,u^*TM)$ 的 $(0,1)$-部分。也就是说，
\[
  (\widetilde\nabla\xi)^{0,1}
  :=\frac12\bigl(
    \widetilde\nabla\xi+J(u)\widetilde\nabla\xi\circ j_\Sigma
  \bigr).
\]
在 $\omega$-\BookRecordIndex{T0372}驯服情形，$D_u$ 的复线性部分由联络
\[
  \widehat\nabla_vX
  :=\widetilde\nabla_vX
  -\frac14(\nabla_{JX}J)v
  -\frac14J(\nabla_XJ)v
\]
确定。这个联络保持 $J$，但未必保持度量；它的\BookRecordIndex{T0248}挠率等于 \BookRecordIndex{T0042}Nijenhuis 张量的负四分之一。
此时算子 $D_u$ 由下式给出：
\begin{equation}\label{eq:3-1-7}
  D_u\xi
  =(\widehat\nabla\xi)^{0,1}
  +\frac14N_J\bigl(\xi,\partial_J(u)\bigr)
  +\frac14\bigl(
    J(\nabla_{\bar\partial_J(u)}J)
    +\nabla_{J\bar\partial_J(u)}J
  \bigr)\xi.
\end{equation}
注意，\eqref{eq:3-1-7} 中后三项是复反线性的，而第一项是复线性的。因此，$D_u$ 分裂为
一个复线性\BookRecordIndex{T0210}柯西--黎曼算子与一个零阶复反线性算子之和。对一般实线性\BookRecordIndex{T0210}柯西--黎曼算子，
这一点仍然成立；见习题 C.1.7。附录 C 证明了每个与 $\omega$ \BookRecordIndex{T0348}相容的 $J$ 都满足
\[
  \nabla_{Jv}J=-J\nabla_vJ,
\]
因而在这种情形下 $\widehat\nabla=\widetilde\nabla$，且 \eqref{eq:3-1-7} 的最后两项
消失。更详细的讨论见引理 C.7.2 和引理 C.7.3。
\end{remark}

\begin{exercise}\label{ex:3-1-4}
证明：对每条 \BookRecordIndex{T0003}$J$-全纯曲线 $u:\Sigma\to M$ 和每个向量场
$X\in\Vect(M)$，都有
\[
  D_u(X\circ u)
  =\frac12(\mathcal L_XJ)(u)du\circ j_\Sigma
  =-\frac12J(u)(\mathcal L_XJ)(u)\partial_J(u).
\]
\textit{提示：}设
$\mathbb R\to\operatorname{Diff}(M):t\mapsto\phi_t$ 是 $X$ 的流，并记
\[
  u_t:=\phi_t^{-1}\circ u,\qquad J_t:=\phi_t^*J.
\]
于是对每个 $t$，$u_t$ 都是一条 $J_t$-全纯曲线。现在关于联络
$\widetilde\nabla$，在 $t=0$ 处对恒等式
$\bar\partial_{J_t}(u_t)=0$ 作协变微分。另一种做法是利用
习题~\ref{ex:2-1-1} 的公式直接计算。
\end{exercise}

到目前为止，我们一直假设 $u$ 光滑。然而，映射 $\mathcal F_u$、从而算子 $D_u$ 的
定义可以推广到 $u$ 属于 $W^{k,p}$ 的情形，其中 $k\ge1$ 是整数，$p>2$ 是实数。
在这种情形，$\mathcal F_u$ 是相应索伯列夫完备化之间的映射：
\[
  \mathcal F_u:
  W^{k,p}(\Sigma,u^*TM)
  \longrightarrow
  W^{k-1,p}(\Sigma,\Lambda^{0,1}\otimes_Ju^*TM).
\]
这个映射与 $J$ 一样光滑，而它在零点的微分 $D_u$ 是同样两个空间之间的有界线性算子。
这里，$W^{k,p}(M,E)$ 表示向量丛 $E\to M$ 的光滑截面空间
$\Omega^0(M,E)$ 关于索伯列夫 $W^{k,p}$-范数的完备化。此外，
$\Lambda^{0,1}:=\Lambda^{0,1}T^*\Sigma$ 表示 $\Sigma$ 上 $(0,1)$ 型 $1$-形式丛。
关于\BookRecordIndex{T0319}索伯列夫空间的更多细节，见附录 B。

\paragraph{主论证纲要。}
公式 \eqref{eq:3-1-4} 表明，$D_u$ 是附录 C 所述的实线性\BookRecordIndex{T0210}柯西--黎曼算子。因此，由
定理 C.1.10，$D_u$ 是\BookRecordIndex{T0128}弗雷德霍姆算子；这意味着它的像是闭的，而核和余核都是有限维的。
这种算子的\BookRecordIndex{T0129}弗雷德霍姆指标定义为核的维数减去余核的维数。由附录 A 的定理 A.1.5，
它在紧（低阶）扰动下保持不变。在当前情形，Riemann--Roch 定理 C.1.10 断言，
\BookRecordIndex{T0129}弗雷德霍姆指标为
\[
  \operatorname{index}D_u
  =n(2-2g)+2c_1(u^*TM),
\]
其中 $g$ 是 $\Sigma$ 的亏格。从代数几何观点看，可如下理解这个公式。在
\eqref{eq:3-1-7} 中已经看到，$D_u$ 的复反线性部分是低阶（紧）扰动，因而可以去掉而
不改变\BookRecordIndex{T0129}弗雷德霍姆指标。$D_u$ 的首项是复线性\BookRecordIndex{T0210}柯西--黎曼算子
\[
  \xi\longmapsto(\widehat\nabla\xi)^{0,1},
\]
所以它在向量丛 $u^*TM$ 上定义一个\BookRecordIndex{T0288}全纯结构。该指标是这个向量丛的 Dolbeault 上同调的
欧拉示性数，由 Riemann--Roch 公式给出。

下面概述一个论证，它证明：对通有的\BookRecordIndex{T0088}殆复结构 $J$，
$\mathcal M^*(A,\Sigma;J)$ 是具有“正确维数”
\[
  \dim\mathcal M^*(A,\Sigma;J)
  =n(2-2g)+2c_1(u^*TM)
\]
的流形。为使这一点成立，必须允许 $J$ 在 $M$ 上某个足够大的\BookRecordIndex{T0088}殆复结构空间
$\mathcal J$ 中变化，以便命题~\ref{prop:3-2-1} 的\BookRecordIndex{T0166}横截性论证能够进行。例如，
$\mathcal J$ 可以是
所有光滑\BookRecordIndex{T0088}殆复结构空间的任意一个在 $C^\infty$ 拓扑下的开子集；并不要求 $M$ 是\BookRecordIndex{T0358}辛流形。
本书语境中最重要的情形是 $M$ 带有\BookRecordIndex{T0363}辛形式 $\omega$，而 $\mathcal J$ 或者是所有
$\omega$-\BookRecordIndex{T0348}相容\BookRecordIndex{T0088}殆复结构的空间，或者是所有 $\omega$-\BookRecordIndex{T0372}驯服\BookRecordIndex{T0088}殆复结构的空间
（见第~\ref{sec:2-1}节）。在所有情形下，$\mathcal J$ 都带通常的 $C^\infty$ 拓扑。
相应的 $C^\ell$ \BookRecordIndex{T0088}殆复结构空间记作 $\mathcal J^\ell$。

基本观察是：当 $\ell$ 足够大时，\BookRecordIndex{T0091}单纯曲线的\BookRecordIndex{T0268}普适模空间
\[
  \mathcal M^*(A,\Sigma;\mathcal J^\ell)
  :=\bigl\{(u,J)\bigm|
  J\in\mathcal J^\ell,\ u\in\mathcal M^*(A,\Sigma;J)\bigr\}
\]
是\BookRecordIndex{T0213}可分\BookRecordIndex{T0053}巴拿赫流形（见命题~\ref{prop:3-2-1}）。此外，投影
\[
  \pi:\mathcal M^*(A,\Sigma;\mathcal J^\ell)
  \longrightarrow\mathcal J^\ell
\]
是弗雷德霍姆映射，因为它在点 $(u,J)$ 处的微分本质上就是算子 $D_u$。特别地，
$\pi$ 与 $D_u$ 有相同的指标，并且 $\pi$ 在 $(u,J)$ 处是满射，当且仅当 $D_u$ 是满射。
因为我们处于\BookRecordIndex{T0053}巴拿赫流形的框架中，可以应用\BookRecordIndex{T0380}隐函数定理（定理 A.3.3）得到：只要 $J$ 是
$\pi$ 的\BookRecordIndex{T0409}正则值，$\mathcal M^*(A,\Sigma;J)$ 确实是一个有限维流形，并且它在 $u$ 处的
切空间就是 $D_u$ 的核。此外，由 萨德--斯梅尔定理（定理 A.5.1），当 $\ell$ 足够大时，
$\pi$ 的\BookRecordIndex{T0409}正则值集在 $\mathcal J^\ell$ 中是（贝尔意义下的）\BookRecordIndex{T0300}剩余集。因此，对
$\mathcal J^\ell$ 的通有元素，$\mathcal M^*(A,\Sigma;J)$ 是具有正确维数的流形。

Taubes 的一个论证使我们能够证明：对光滑\BookRecordIndex{T0088}殆复结构空间 $\mathcal J$ 的通有元素，这个
结论仍然成立。该论证基于如下观察：所有单纯 \BookRecordIndex{T0003}$J$-全纯曲线组成的空间是可数个紧集之并，
而对每个这样的紧集，相应的正则 $J$ 之集都是开的。因此，所有正则 $J$ 之集是可数个
开集 $\mathcal J_{\mathrm{reg},K}$ 的交；问题的实质在于证明这些开集也是稠密的。
上面概述的论证通过考虑 $C^\ell$ \BookRecordIndex{T0088}殆复结构的\BookRecordIndex{T0053}巴拿赫流形来做到这一点。
在注~\ref{rem:3-2-7} 中，我们将说明 Floer 提出的另一种办法；它始终留在
$C^\infty$ 范畴内。

以上断言将在第~\ref{sec:3-2}节证明。这里先归纳主要结论。首先给出一个定义。

\begin{definition}\label{def:3-1-5}
固定紧黎曼曲面 $(\Sigma,j_\Sigma,d\operatorname{vol}_\Sigma)$ 和同调类
$A\in H_2(M;\mathbb Z)$。如果对每个
$u\in\mathcal M^*(A,\Sigma;J)$，$D_u$ 都是满射，就称 $M$ 上的\BookRecordIndex{T0088}殆复结构 $J$
（关于 $A$ 和 $\Sigma$）是\emph{正则的}。给定如上的 $\mathcal J$，以
\[
  \mathcal J_{\mathrm{reg}}(A,\Sigma)
\]
表示 $\mathcal J$ 中所有关于 $A$ 和 $\Sigma$ 正则的 $J$ 组成的集合。在
$\Sigma=S^2$ 的情形，简记
\[
  \mathcal J_{\mathrm{reg}}(A)
  :=\mathcal J_{\mathrm{reg}}(A,S^2).
\]
\end{definition}

\begin{theorem}\label{thm:3-1-6}
假设 $\mathcal J=\mathcal J(M,\omega)$ 或
$\mathcal J=\mathcal J_\tau(M,\omega)$。固定紧黎曼曲面
$(\Sigma,j_\Sigma,d\operatorname{vol}_\Sigma)$ 和同调类
$A\in H_2(M;\mathbb Z)$。
\begin{enumerate}[label=(\roman*),leftmargin=3em]
  \item 若 $J\in\mathcal J_{\mathrm{reg}}(A,\Sigma)$，则空间
  $\mathcal M^*(A,\Sigma;J)$ 是维数为
  \[
    \dim\mathcal M^*(A,\Sigma;J)
    =n(2-2g)+2c_1(A)
  \]
  的光滑流形，并带有自然定向。

  \item 集合 $\mathcal J_{\mathrm{reg}}(A,\Sigma)$ 在 $\mathcal J$ 中是剩余的，
  即它包含 $\mathcal J$ 中可数个开稠密子集的交。
\end{enumerate}
\end{theorem}

下一项任务是讨论流形 $\mathcal M^*(A,\Sigma;J)$ 对
$J\in\mathcal J_{\mathrm{reg}}(A,\Sigma)$ 之选取的依赖。\BookRecordIndex{T0088}殆复结构的一个（光滑）同伦，
是光滑映射
\[
  [0,1]\longrightarrow\mathcal J:
  \lambda\longmapsto J_\lambda.
\]
对任意这种同伦，定义
\[
  \mathcal W^*(A,\Sigma;\{J_\lambda\}_\lambda)
  :=\bigl\{(\lambda,u)\bigm|
  0\le\lambda\le1,\ u\in\mathcal M^*(A,\Sigma;J_\lambda)\bigr\}.
\]
给定 $J_0,J_1\in\mathcal J_{\mathrm{reg}}(A,\Sigma)$，以
$\mathcal J(J_0,J_1)$ 表示连接 $J_0$ 与 $J_1$ 的所有光滑\BookRecordIndex{T0088}殆复结构同伦组成的空间。
一般而言，尽管 $\mathcal J$ 是道路连通的，却不一定存在一个对每个 $\lambda$ 都满足
$J_\lambda\in\mathcal J_{\mathrm{reg}}(A,\Sigma)$ 的同伦。换言之，对某些 $\lambda$，
空间 $\mathcal M^*(A,\Sigma;J_\lambda)$ 可能不是流形。不过，我们将看到，总能找到一个
光滑同伦，使 $\mathcal W^*(A,\Sigma;\{J_\lambda\}_\lambda)$ 成为流形。

\begin{definition}\label{def:3-1-7}
固定紧黎曼曲面 $(\Sigma,j_\Sigma,d\operatorname{vol}_\Sigma)$ 和同调类
$A\in H_2(M;\mathbb Z)$。设
$J_0,J_1\in\mathcal J_{\mathrm{reg}}(A,\Sigma)$。若从 $J_0$ 到 $J_1$ 的同伦
\[
  [0,1]\longrightarrow\mathcal J:
  \lambda\longmapsto J_\lambda
\]
对每个 $(\lambda,u)\in\mathcal W^*(A,\Sigma;\{J_\lambda\}_\lambda)$ 都满足
\[
  \Omega^{0,1}(\Sigma,u^*TM)
  =\operatorname{im}D_{J_\lambda,u}+\mathbb Rv_\lambda,
\]
其中
\[
  v_\lambda:=(\partial_\lambda J_\lambda)du\circ j_\Sigma
\]
是道路 $\lambda\mapsto J_\lambda$ 的切向量在
$\Omega^{0,1}(\Sigma,u^*TM)$ 中的像，则称该同伦（关于 $A$ 和 $\Sigma$）
是\emph{正则的}。\BookRecordIndex{T0407}正则同伦的空间记作
$\mathcal J_{\mathrm{reg}}(A,\Sigma;J_0,J_1)$。
\end{definition}

我们将在第~\ref{sec:3-2}节看到，同伦 $\lambda\mapsto J_\lambda$ 正则，当且仅当它与投影
$\pi$ 横截。直观地说，可以把正则\BookRecordIndex{T0088}殆复结构空间
$\mathcal J_{\mathrm{reg}}(A,\Sigma)$ 想成 $\mathcal J$ 中一个余维为一的子簇
$\mathcal S$ 的补集。光滑同伦 $\lambda\mapsto J_\lambda$ 正则，恰好在它与
$\mathcal S$ 横截时成立。

\begin{theorem}\label{thm:3-1-8}
假设 $\mathcal J=\mathcal J(M,\omega)$ 或
$\mathcal J=\mathcal J_\tau(M,\omega)$。固定紧黎曼曲面
$(\Sigma,j_\Sigma,d\operatorname{vol}_\Sigma)$ 和同调类
$A\in H_2(M;\mathbb Z)$。设
$J_0,J_1\in\mathcal J_{\mathrm{reg}}(A,\Sigma)$。
\begin{enumerate}[label=(\roman*),leftmargin=3em]
  \item 若
  $\{J_\lambda\}_\lambda\in
  \mathcal J_{\mathrm{reg}}(A,\Sigma;J_0,J_1)$，则
  $\mathcal W^*(A,\Sigma;\{J_\lambda\}_\lambda)$ 是光滑定向流形，其边界为
  \[
    \partial\mathcal W^*(A,\Sigma;\{J_\lambda\}_\lambda)
    =\mathcal M^*(A,\Sigma;J_0)\cup
    \mathcal M^*(A,\Sigma;J_1).
  \]
  边界定向与 $\mathcal M^*(A,\Sigma;J_1)$ 的定向一致，而与
  $\mathcal M^*(A,\Sigma;J_0)$ 的定向相反。

  \item 集合
  $\mathcal J_{\mathrm{reg}}(A,\Sigma;J_0,J_1)$
  在所有从 $J_0$ 到 $J_1$ 的光滑同伦
  $[0,1]\to\mathcal J:\lambda\mapsto J_\lambda$ 所成的空间中是剩余的。
\end{enumerate}
\end{theorem}

定理~\ref{thm:3-1-8} 表明，\BookRecordIndex{T0246}模空间
$\mathcal M^*(A,\Sigma;J_0)$ 与
$\mathcal M^*(A,\Sigma;J_1)$ \BookRecordIndex{T0107}定向配边。不过，在建立某种紧致性之前，这并没有太多
意义。这个问题将在第 4 章处理。我们用来建立\BookRecordIndex{T0166}横截性和紧致性的论证，都依赖下面所解释的
\BookRecordIndex{T0329}椭圆正则性基本观察，但此后两种论证便分道扬镳，实质上彼此独立。尽管如此，为理解我们的
记号并对\BookRecordIndex{T0319}索伯列夫空间有所认识，仍建议不熟悉这些主题的读者继续读到下一节末尾，然后再
转去考虑紧致性问题。

\begin{remark}\label{rem:3-1-9}
这里也许正适合提醒一句。在我们的讨论中，“正则”一词有两种相当不同的用法。第一种
（如“\BookRecordIndex{T0329}椭圆正则性”）谈的是映射的光滑程度，例如一个给定映射 $u:\Sigma\to M$
究竟属于哪个 $W^{k,p}$ 函数空间。另一种（如“正则的 $J$”）所关心的是 $J$ 是否为
投影映射 $\pi$ 的\BookRecordIndex{T0409}正则值而不是\BookRecordIndex{T0234}临界值。人们常用“通有的”代替“正则的”来修饰 $J$。
这种术语的困难在于，“通有”一词还带有别的含义；例如，说 $S^2\times S^2$ 上的乘积
\BookRecordIndex{T0088}殆复结构是通有的，似乎自相矛盾，因为它如此特殊。另一方面，它对于所有亏格为零的曲线
都是正则的。
\end{remark}

\paragraph{\BookRecordIndex{T0329}椭圆正则性。}
为准备定理~\ref{thm:3-1-6} 与定理~\ref{thm:3-1-8} 的证明，我们先引入适当的
\BookRecordIndex{T0319}索伯列夫空间框架，并讨论相关的\BookRecordIndex{T0329}椭圆正则性定理。证明见附录 B 和附录 C。

此后，为了得到具有\BookRecordIndex{T0053}巴拿赫流形结构的参数空间，考虑 $C^\ell$ 而非光滑的\BookRecordIndex{T0088}殆复结构会很
重要。因此，我们将考虑所有 $C^\ell$ 的 $\omega$-\BookRecordIndex{T0372}驯服\BookRecordIndex{T0088}殆复结构所成的空间
\[
  \mathcal J^\ell=\mathcal J_\tau^\ell(M,\omega).
\]
下文假设 $\ell\ge2$，从而由 $\omega$ 和 $J$ 确定的度量 \eqref{eq:2-1-1} 至少属于
$C^2$，第~\ref{sec:2-5}节的结果也适用。不过，本节讨论的\BookRecordIndex{T0408}正则性结果对 $\ell=1$
仍然成立。

对实数 $p>2$ 和整数 $k\ge1$，记
\[
  \mathcal B^{k,p}
  :=\bigl\{u\in W^{k,p}(\Sigma,M)\bigm|[u]=A\bigr\},
\]
即由代表类 $A\in H_2(M;\mathbb Z)$、且 $k$ 阶导数属于 $L^p$ 的连续映射
$\Sigma\to M$ 组成的空间。如附录 B 更详细说明的，空间 $W^{k,p}(\Sigma,M)$ 是
$C^\infty(\Sigma,M)$ 关于一个以索伯列夫 $W^{k,p}$-范数为基础的距离函数的完备化。
这个范数定义为所有不超过 $k$ 阶的导数的 $L^p$-范数之和。可以把 $M$ 嵌入某个欧氏
空间 $\mathbb R^N$，再利用环境空间 $W^{k,p}(\Sigma,\mathbb R^N)$ 上的索伯列夫范数，
来定义 $C^\infty(\Sigma,M)$ 上相应的度量。另一种做法是，在 $M$ 上选取黎曼度量，
用协变导数定义切空间
\[
  T_uC^\infty(\Sigma,M)=\Omega^0(\Sigma,u^*TM)
\]
上的 $W^{k,p}$-范数，再对 $C^\infty(\Sigma,M)$ 中端点固定的道路长度取下确界。
由于 $\Sigma$ 与 $M$ 都紧致，任意两个这样的 $W^{k,p}$ 距离函数都彼此\BookRecordIndex{T0348}相容，所以所得
完备化与这些选取无关。完备化 $W^{k,p}(\Sigma,M)$ 是光滑\BookRecordIndex{T0213}可分\BookRecordIndex{T0053}巴拿赫流形，其切空间为
\[
  T_uW^{k,p}(\Sigma,M)=W^{k,p}(\Sigma,u^*TM).
\]
局部坐标卡可以利用沿光滑映射 $u:\Sigma\to M$ 的指数映射来定义；见下面
命题~\ref{prop:3-2-1} 的证明。空间 $\mathcal B^{k,p}$ 是
$W^{k,p}(\Sigma,M)$ 的若干连通分支之并
\sourcefootnote{原书称 $\mathcal B^{k,p}$ 是 $W^{k,p}(\Sigma,M)$ 的一个连通分支；
但固定同调类一般不能唯一确定映射空间的同伦分支，故改作“若干连通分支之并”。}，
因而它也是光滑\BookRecordIndex{T0213}可分\BookRecordIndex{T0053}巴拿赫流形。

为了使空间 $\mathcal B^{k,p}$ 定义良好，必须假设
\[
  kp>2.
\]
这有几个原因。第一，用 $W^{k,p}$ 类的局部坐标表示来定义 $\mathcal B^{k,p}$ 本身就
需要这个假设。要点如下：对欧氏空间中的开集之间的映射，$W^{k,p}$-范数是良定义的；
但在一般流形上，需要一个在与坐标卡复合下不变的空间。一个
$W^{k,p}$ 映射 $u:\mathbb R^2\to\mathbb R^{2n}$ 与靶空间中的 $C^k$ 映射复合后
仍属于 $W^{k,p}$，恰好需要 $kp>2$。第二，我们将经常使用\BookRecordIndex{T0318}索伯列夫嵌入定理；它说在
这个条件下，空间 $\mathcal B^{k,p}$ 嵌入 $\Sigma\to M$ 的连续映射空间。此外，
由 Rellich 定理，这个嵌入是紧的。第三，要证明两个 $W^{k,p}$ 函数的乘积仍然属于
$W^{k,p}$，也需要条件 $kp>2$。换言之，处理非线性项需要这个条件。

第一个关键观察是：当 $J$ 光滑时，每条 $W^{1,p}$ 类的 \BookRecordIndex{T0003}$J$-全纯曲线（其中 $p>2$）
必然光滑。更准确地说，有下面这个可由\BookRecordIndex{T0330}椭圆自举法证明的\BookRecordIndex{T0408}正则性定理。细节见附录 B
（定理 B.4.1 和注 B.4.3）。

\begin{proposition}[\BookRecordIndex{T0329}椭圆正则性]\label{prop:3-1-10}
假设 $J$ 是一个 $C^\ell$ \BookRecordIndex{T0088}殆复结构，其中 $\ell\ge1$。若
$u:\Sigma\to M$ 是一条 $W^{1,p}$ 类的 \BookRecordIndex{T0003}$J$-全纯曲线，其中 $p>2$，则
$u$ 属于 $W^{\ell+1,p}$。特别地，$u$ 属于 $C^\ell$；若 $J$ 光滑
（即属于 $C^\infty$），则 $u$ 也光滑。
\end{proposition}

命题~\ref{prop:3-1-10} 表明：对 $J\in\mathcal J^\ell$，由 $W^{k,p}$ 类
\BookRecordIndex{T0003}$J$-全纯曲线组成的\BookRecordIndex{T0246}模空间与 $k$ 的选取无关，只要 $k\le\ell+1$。方程线性化后会损失
$J$ 的一阶导数，所以需要条件 $k\le\ell$，才能在相应的\BookRecordIndex{T0319}索伯列夫空间上良定义算子
$D_u$。下一个结果是命题~\ref{prop:3-1-10} 的线性版本。它把算子
\begin{equation}\label{eq:3-1-8}
  D_u:W^{k,p}(\Sigma,u^*TM)
  \longrightarrow
  W^{k-1,p}(\Sigma,\Lambda^{0,1}\otimes_Ju^*TM)
\end{equation}
的余核与\BookRecordIndex{T0368}形式伴随算子
\begin{equation}\label{eq:3-1-9}
  D_u^*:W^{k,p}(\Sigma,\Lambda^{0,1}\otimes_Ju^*TM)
  \longrightarrow
  W^{k-1,p}(\Sigma,u^*TM)
\end{equation}
的核联系起来，反之亦然。算子 $D_u^*$ 由恒等式
\[
  \int_\Sigma\langle\eta,D_u\xi\rangle\,d\operatorname{vol}_\Sigma
  =
  \int_\Sigma\langle D_u^*\eta,\xi\rangle\,d\operatorname{vol}_\Sigma
\]
定义，其中 $\xi\in W^{k,p}(\Sigma,u^*TM)$，且
$\eta\in W^{k,p}(\Sigma,\Lambda^{0,1}\otimes_Ju^*TM)$。它是一阶微分算子，也可解释为
实线性\BookRecordIndex{T0210}柯西--黎曼算子（见附录 C）。

\begin{proposition}\label{prop:3-1-11}
固定正整数 $\ell$ 和常数 $p>2$。设 $J$ 是 $M$ 上一个 $\omega$-\BookRecordIndex{T0372}驯服的
$C^\ell$ \BookRecordIndex{T0088}殆复结构，并设 $u\in W^{\ell,p}(\Sigma,M)$。令
$k\in\{1,\ldots,\ell\}$，并令 $q>1$ 满足 $1/p+1/q=1$。则有以下结论。
\begin{enumerate}[label=(\roman*),leftmargin=3em]
  \item 算子 \eqref{eq:3-1-8} 与 \eqref{eq:3-1-9} 都是\BookRecordIndex{T0128}弗雷德霍姆算子，且
  \[
    \operatorname{index}D_u
    =-\operatorname{index}D_u^*
    =n(2-2g)+2c_1(u^*TM),
  \]
  其中 $g$ 是 $\Sigma$ 的亏格。

  \item 若
  $\eta\in L^q(\Sigma,\Lambda^{0,1}\otimes_Ju^*TM)$ 满足
  \[
    \int_\Sigma\langle\eta,D_u\xi\rangle
    \,d\operatorname{vol}_\Sigma=0
  \]
  对所有 $\xi\in W^{k,p}(\Sigma,u^*TM)$ 都成立，则
  $\eta\in W^{\ell,p}(\Sigma,\Lambda^{0,1}\otimes_Ju^*TM)$，并且
  $D_u^*\eta=0$。

  \item 若 $\xi\in L^q(\Sigma,u^*TM)$ 满足
  \[
    \int_\Sigma\langle\xi,D_u^*\eta\rangle
    \,d\operatorname{vol}_\Sigma=0
  \]
  对所有
  $\eta\in W^{k,p}(\Sigma,\Lambda^{0,1}\otimes_Ju^*TM)$ 都成立，则
  $\xi\in W^{\ell,p}(\Sigma,u^*TM)$，并且 $D_u\xi=0$。
\end{enumerate}
\end{proposition}

\begin{proof}
先假设 $u$ 光滑，并考虑 $u^*TM$ 上的联络
\[
  \Omega^0(\Sigma,u^*TM)\longrightarrow\Omega^1(\Sigma,u^*TM):
  \xi\longmapsto
  \nabla\xi-\frac12J(u)(\nabla_\xi J)(u)du.
\]
一般而言，这个联络既不保持复结构，也不保持度量。
算子 $D_u$ 是这个映射与投影
\[
  \Omega^1(\Sigma,u^*TM)
  \longrightarrow\Omega^{0,1}(\Sigma,u^*TM)
\]
的复合，并延拓到相应的索伯列夫完备化。因此，$D_u$ 是实线性\BookRecordIndex{T0210}柯西--黎曼算子。
若 $u$ 属于 $W^{\ell,p}$，同一论证表明，$D_u$ 是定义 C.1.5 意义下
$W^{\ell-1,p}$ 类的实线性\BookRecordIndex{T0210}柯西--黎曼算子。因此，结论由定理 C.2.3 得出；在该定理的
记号中取 $q=p$，并令 $1/r+1/p=1$。
\end{proof}

命题~\ref{prop:3-1-11} 蕴含，$D_u$ 的核与余核不依赖于定义该算子时对空间的确切选择。
为看出这一点，假设 $J$ 属于 $C^\ell$，且 $u$ 是一条 \BookRecordIndex{T0003}$J$-全纯曲线。于是，由
命题~\ref{prop:3-1-10}，$u$ 属于 $C^\ell$。现在考虑 $k\le\ell$ 时的算子
\eqref{eq:3-1-8}。由命题~\ref{prop:3-1-11}，$D_u$ 核中的每个元素对任意 $p$
都必属于 $W^{\ell,p}$，所以只要 $k\le\ell$，$D_u$ 的核就与 $k$ 和 $p$ 的选取无关。
余核也有同样的结论。特别地，$D_u$ 对某个 $k,p$ 的选择是满射，当且仅当它对所有这些
选择都是满射。
